%**************************************%
%*    Generated from PreTeXt source   *%
%*    on 2018-08-18T17:55:39Z    *%
%*                                    *%
%*   http://mathbook.pugetsound.edu   *%
%*                                    *%
%**************************************%
\documentclass[10pt,]{article}
%% Custom Preamble Entries, early (use latex.preamble.early)
%% Default LaTeX packages
%%   1.  always employed (or nearly so) for some purpose, or
%%   2.  a stylewriter may assume their presence
\usepackage{geometry}
%% Some aspects of the preamble are conditional,
%% the LaTeX engine is one such determinant
\usepackage{ifthen}
%% etoolbox has a variety of modern conveniences
\usepackage{etoolbox}
\usepackage{ifxetex,ifluatex}
%% Raster graphics inclusion
\usepackage{graphicx}
%% Color support, xcolor package
%% Always loaded, for: add/delete text, author tools
%% Here, since tcolorbox loads tikz, and tikz loads xcolor
\PassOptionsToPackage{usenames,dvipsnames,svgnames,table}{xcolor}
\usepackage{xcolor}
%% Colored boxes, and much more, though mostly styling
%% skins library provides "enhanced" skin, employing tikzpicture
%% boxes may be configured as "breakable" or "unbreakable"
%% "raster" controls grids of boxes, aka side-by-side
\usepackage{tcolorbox}
\tcbuselibrary{skins}
\tcbuselibrary{breakable}
\tcbuselibrary{raster}
%% xparse allows the construction of more robust commands,
%% this is a necessity for isolating styling and behavior
%% The tcolorbox library of the same name loads the base library
\tcbuselibrary{xparse}
%% Hyperref should be here, but likes to be loaded late
%%
%% Inline math delimiters, \(, \), need to be robust
%% 2016-01-31:  latexrelease.sty  supersedes  fixltx2e.sty
%% If  latexrelease.sty  exists, bugfix is in kernel
%% If not, bugfix is in  fixltx2e.sty
%% See:  https://tug.org/TUGboat/tb36-3/tb114ltnews22.pdf
%% and read "Fewer fragile commands" in distribution's  latexchanges.pdf
\IfFileExists{latexrelease.sty}{}{\usepackage{fixltx2e}}
%% Text height identically 9 inches, text width varies on point size
%% See Bringhurst 2.1.1 on measure for recommendations
%% 75 characters per line (count spaces, punctuation) is target
%% which is the upper limit of Bringhurst's recommendations
\geometry{letterpaper,total={340pt,9.0in}}
%% Custom Page Layout Adjustments (use latex.geometry)
%% This LaTeX file may be compiled with pdflatex, xelatex, or lualatex
%% The following provides engine-specific capabilities
%% Generally, xelatex and lualatex will do better languages other than US English
%% You can pick from the conditional if you will only ever use one engine
\ifthenelse{\boolean{xetex} \or \boolean{luatex}}{%
%% begin: xelatex and lualatex-specific configuration
%% fontspec package will make Latin Modern (lmodern) the default font
\ifxetex\usepackage{xltxtra}\fi
\usepackage{fontspec}
%% realscripts is the only part of xltxtra relevant to lualatex 
\ifluatex\usepackage{realscripts}\fi
%% 
%% Extensive support for other languages
\usepackage{polyglossia}
%% Main document language is US English
\setdefaultlanguage{english}
%% Spanish
\setotherlanguage{spanish}
%% Vietnamese
\setotherlanguage{vietnamese}
%% end: xelatex and lualatex-specific configuration
}{%
%% begin: pdflatex-specific configuration
%% translate common Unicode to their LaTeX equivalents
%% Also, fontenc with T1 makes CM-Super the default font
%% (\input{ix-utf8enc.dfu} from the "inputenx" package is possible addition (broken?)
\usepackage[T1]{fontenc}
\usepackage[utf8]{inputenc}
%% end: pdflatex-specific configuration
}
%% Symbols, align environment, bracket-matrix
\usepackage{amsmath}
\usepackage{amssymb}
%% allow page breaks within display mathematics anywhere
%% level 4 is maximally permissive
%% this is exactly the opposite of AMSmath package philosophy
%% there are per-display, and per-equation options to control this
%% split, aligned, gathered, and alignedat are not affected
\allowdisplaybreaks[4]
%% allow more columns to a matrix
%% can make this even bigger by overriding with  latex.preamble.late  processing option
\setcounter{MaxMatrixCols}{30}
%%
%% Semantic Macros
%% To preserve meaning in a LaTeX file
%% Only defined here if required in this document
%% Used for inline definitions of terms
\newcommand{\terminology}[1]{\textbf{#1}}
%% Subdivision Numbering, Chapters, Sections, Subsections, etc
%% Subdivision numbers may be turned off at some level ("depth")
%% A section *always* has depth 1, contrary to us counting from the document root
%% The latex default is 3.  If a larger number is present here, then
%% removing this command may make some cross-references ambiguous
%% The precursor variable $numbering-maxlevel is checked for consistency in the common XSL file
\setcounter{secnumdepth}{3}
%% begin: General AMS environment setup
%% Environments built with amsthm package
\usepackage{amsthm}
%% Numbering for Theorems, Conjectures, Examples, Figures, etc
%% Controlled by  numbering.theorems.level  processing parameter
%% Numbering: all theorem-like numbered consecutively
%% i.e. Corollary 4.3 follows Theorem 4.2
%% Always need some theorem environment to set base numbering scheme
%% even if document has no theorems (but has other environments)
%% Create a never-used style first, always
%% simply to provide a global counter to use, namely "cthm"
\newtheorem{cthm}{BadTheoremStringName}[section]
%% end: General AMS environment setup
%% begin: environments with italicized bodies, theorems and similar
%% Style is like a theorem, and for statements without proofs
%% Theorem-like environments in modified "plain" style
%% We manage the head, do not adjust vertical spacing
%% Thus "space after theorem head" is necessary
%% This provides an automatic period after the number
\newtheoremstyle{ptxplainnotitle}
  {}% space above
  {}% space below
  {\itshape}% body font
  {}% indent amount
  {\bfseries}% theorem head font
  {.}% punctuation after theorem head
  {0.5em}% space after theorem head
  {}% theorem head specification
%% We now manage punctuation on-sight, elsewhere,
%% assuming non-trivial content inside a "title"
\newtheoremstyle{ptxplaintitle}
  {}% space above
  {}% space below
  {\itshape}% body font
  {}% indent amount
  {\bfseries}% theorem head font
  {}% punctuation after theorem head
  {0.5em}% space after theorem head
  {\thmname{#1}\thmnumber{ #2}\thmnote{ #3}}% theorem head specification
%% Only variants actually used in document appear here
%% Template eventually creates an environment of the given name
%% No arguments => Theorem 2.6. via "notitle" style variant
%% One optional argument => Theorem 2.6 Fantastic! via "title" style variant
%% AMS proof environment is basically fine as-is and special treatment
%% would certainly interfere with the functioning of \qed, etc.
%% So we simply localize the default heading
%% Redefinition of the "proof" environment is to cause a long alternate
%% title to line-break appropriately.  Code is cut verbatim, by suggestion,
%% from "Using the amsthm Package" Version 2.20.3, September 2017
%% end: environments with italicized bodies, theorems and similar
%% begin: environments with normal bodies, examples, etc.
%% Other environments go in modified "definition" style
%% Similar to above
\newtheoremstyle{ptxdefinitionnotitle}
  {}% space above
  {}% space below
  {}% body font
  {}% indent amount
  {\bfseries}% theorem head font
  {.}% punctuation after theorem head
  {0.5em}% space after theorem head
  {\thmname{#1}\thmnumber{ #2}}% theorem head specification
\newtheoremstyle{ptxdefinitiontitle}
  {}% space above
  {}% space below
  {}% body font
  {}% indent amount
  {\bfseries}% theorem head font
  {}% punctuation after theorem head
  {0.5em}% space after theorem head
  {\thmname{#1}\thmnumber{ #2}\thmnote{ #3}}% theorem head specification
%% definition: fairly simple numbered block/structure
\tcbset{ definitionstyle/.style={size=minimal, colback=white, colbacktitle=white, coltitle=black, fonttitle=\bfseries, attach title to upper, after title={\space}} }
\newtcolorbox[use counter=cthm, number within=section]{definition}[2]{title={{Definition~\thetcbcounter\notblank{#1}{\space\space#1}{}}}, label=#2, breakable, definitionstyle}
%% end: environments with normal bodies, examples, etc.
%% Localize LaTeX supplied names (possibly none)
\renewcommand*{\appendixname}{Appendix}
\renewcommand*{\abstractname}{Abstract}
%% Figures, Tables, Listings, Named Lists, Floats
%% The [H]ere option of the float package fixes floats in-place,
%% in deference to web usage, where floats are totally irrelevant
%% You can remove some of this setup, to restore standard LaTeX behavior
%% HOWEVER, numbering of figures/tables AND theorems/examples/remarks, etc
%% may de-synchronize with the numbering in the HTML version
%% You can remove the "placement={H}" option to allow flotation and
%% preserve numbering, BUT the numbering may then appear "out-of-order"
%% Floating environments: http://tex.stackexchange.com/questions/95631/
\usepackage{float}
\usepackage{newfloat}
\usepackage{caption}%% Adjust stock figure environment so that it no longer floats
\SetupFloatingEnvironment{figure}{fileext=lof,placement={H},within=section,name=Figure}
\captionsetup[figure]{labelfont=bf}
%% http://tex.stackexchange.com/questions/16195
\makeatletter
\let\c@figure\c@cthm
\makeatother
%% Footnote Numbering
%% We reset the footnote counter, as given by numbering.footnotes.level
\makeatletter\@addtoreset{footnote}{section}\makeatother
%% More flexible list management, esp. for references
%% But also for specifying labels (i.e. custom order) on nested lists
\usepackage{enumitem}
%% hyperref driver does not need to be specified, it will be detected
\usepackage{hyperref}
%% configure hyperref's  \url  to match listings' inline verbatim
\renewcommand\UrlFont{\small\ttfamily}
%% Hyperlinking active in PDFs, all links solid and blue
\hypersetup{colorlinks=true,linkcolor=blue,citecolor=blue,filecolor=blue,urlcolor=blue}
\hypersetup{pdftitle={Intro to Topology}}
%% If you manually remove hyperref, leave in this next command
\providecommand\phantomsection{}
%% If tikz has been loaded, replace ampersand with \amp macro
%% extpfeil package for certain extensible arrows,
%% as also provided by MathJax extension of the same name
%% NB: this package loads mtools, which loads calc, which redefines
%%     \setlength, so it can be removed if it seems to be in the 
%%     way and your math does not use:
%%     
%%     \xtwoheadrightarrow, \xtwoheadleftarrow, \xmapsto, \xlongequal, \xtofrom
%%     
%%     we have had to be extra careful with variable thickness
%%     lines in tables, and so also load this package late
\usepackage{extpfeil}
%% Custom Preamble Entries, late (use latex.preamble.late)
%% Begin: Author-provided packages
%% (From  docinfo/latex-preamble/package  elements)
%% End: Author-provided packages
%% Begin: Author-provided macros
%% (From  docinfo/macros  element)
%% Plus three from MBX for XML characters
\newcommand{\tuple}[1]{\langle #1 \rangle}
\newcommand{\mb}{\mathbb}
\newcommand{\lt}{<}
\newcommand{\gt}{>}
\newcommand{\amp}{&}
%% End: Author-provided macros
%% Title page information for article
\title{Intro to Topology}
\author{Steven Clontz\\
University of South Alabama
}
\date{August 18, 2018}
\begin{document}
%% Target for xref to top-level element is document start
\hypertarget{index}{}
\maketitle
\thispagestyle{empty}
\begin{abstract}
\hypertarget{p-1}{}%
These are course notes for an introductory course for general topology, to be used in an inquiry-based learning classroom.%
\par
\hypertarget{p-2}{}%
Much of the content has been adapted from Michel Smith's undergraduate topology notes, used with permission.%
\end{abstract}
\hypertarget{p-3}{}%
These are course notes for an introductory course for general topology, to be used in an inquiry-based learning classroom.%
\typeout{************************************************}
\typeout{Section 1 Curves and Surfaces}
\typeout{************************************************}
\section[{Curves and Surfaces}]{Curves and Surfaces}\label{section-surfaces}
\hypertarget{p-4}{}%
In this section, we will develop an intuition for a topological space and the purpose of topology by investigating two natural examples of topological spaces: curves and surfaces.%
\par
\hypertarget{p-5}{}%
Unlike the rest of these notes, this section is not meant to be studied rigorously. (For example, what do I mean by ``locally looks like'' in \hyperref[curve-defn]{Definition~\ref{curve-defn}}?)  However, many of these ideas will return later in the course and be handled more carefully.%
\begin{definition}{}{curve-defn}%
\hypertarget{p-6}{}%
A \terminology{curve} is a set of points such that for every point in the set, the set locally looks like a (possibly bent or curved) copy of the real line \(\mathbb R\) or the half line \(\mathbb R^*=\{x\in\mathbb R:x\geq 0\}\).%
\end{definition}
\hypertarget{p-7}{}%
For example, \hyperref[circle-fig]{Figure~\ref{circle-fig}}, \hyperref[figure-parabola]{Figure~\ref{figure-parabola}}, and \hyperref[figure-helix]{Figure~\ref{figure-helix}} are all examples of curves in two or three-dimensional Euclidean space.%
\begin{figure}
\centering
\includegraphics[width=1\linewidth]{images/math.jpg}
\caption{A circle\label{circle-fig}}
\end{figure}
\begin{figure}
\centering
\includegraphics[width=1\linewidth]{images/math.jpg}
\caption{A parabola\label{figure-parabola}}
\end{figure}
\begin{figure}
\centering
\includegraphics[width=1\linewidth]{images/math.jpg}
\caption{A helix\label{figure-helix}}
\end{figure}
\hypertarget{p-8}{}%
Note the following differences between \hyperref[circle-fig]{Figure~\ref{circle-fig}} and \hyperref[figure-parabola]{Figure~\ref{figure-parabola}}:%
\leavevmode%
\begin{itemize}[label=\textbullet]
\item{}Removing a point from \hyperref[figure-parabola]{Figure~\ref{figure-parabola}} would split it into two disconnected parts, but \hyperref[circle-fig]{Figure~\ref{circle-fig}} would remain connected after a point is removed.%
\item{}\hyperref[circle-fig]{Figure~\ref{circle-fig}} is bounded while \hyperref[figure-parabola]{Figure~\ref{figure-parabola}} extends unboundedly. \footnote{This topological distinction makes sense as both are closed subsets of \(\mathbb R^2\); see \hyperref[section-compactness]{Section~\ref{section-compactness}} for more info.\label{fn-1}}%
\end{itemize}
\hypertarget{p-9}{}%
These differences would remain no matter how the curves were stretched or bent. However, while there are certainly geometrical differences beteween \hyperref[figure-parabola]{Figure~\ref{figure-parabola}} and \hyperref[figure-helix]{Figure~\ref{figure-helix}}, they are in a certain sense the same object that has been bent or stretched into a different shape.%
\begin{definition}{}{definition-2}%
\hypertarget{p-10}{}%
Two objects are said to be \terminology{topologically equivalent} or \terminology{homeomorphic} if one may be bent or stretched into the shape of the other.%
\end{definition}
\hypertarget{p-11}{}%
So this means that all geometrically similar shapes are homeomorphic (as in \hyperref[figure-geometrically-similar]{Figure~\ref{figure-geometrically-similar}}), but we also use the idea of homeomorphism to compare other objects in our daily lives.%
\par
\hypertarget{p-12}{}%
For example, while many of them are not curves by our definition, the letters of the alphabet may be considered as topological objects. \hyperref[figure-letter-a]{Figure~\ref{figure-letter-a}} illustrates several homeomorphic expressions of the letter ``A''.%
\begin{figure}
\centering
\includegraphics[width=1\linewidth]{images/math.jpg}
\caption{Two similar triangles\label{figure-geometrically-similar}}
\end{figure}
\begin{figure}
\centering
\includegraphics[width=1\linewidth]{images/math.jpg}
\caption{The letter ``A'' in several fonts.\label{figure-letter-a}}
\end{figure}
\hypertarget{p-13}{}%
A homeomorphism is more carefully defined in \hyperref[section-continuity]{Section~\ref{section-continuity}}, but the central idea is that of ``neighborhoods''. For each of the letters ``A'' in \hyperref[figure-letter-a]{Figure~\ref{figure-letter-a}}, note that there are two endpoints and two triad intersections whose neighborhoods look different from the other neighborhoods within the letter; see \hyperref[figure-letter-a-neighborhoods]{Figure~\ref{figure-letter-a-neighborhoods}}.%
\begin{figure}
\centering
\includegraphics[width=1\linewidth]{images/math.jpg}
\caption{Neighborhoods within the letter ``A''.\label{figure-letter-a-neighborhoods}}
\end{figure}
\begin{definition}{}{definition-surface}%
\hypertarget{p-14}{}%
A \terminology{surface} is a set of points such that for every point in the set, the set locally looks like a (possibly bent or curved) copy of the plane \(\mathbb R^2\) or the half-plane \(\mathbb R^{2*}=\{\tuple{x,y}\in\mb R^2:x\geq 0\}\).%
\end{definition}
\hypertarget{p-15}{}%
A classic example of the topology of surfaces is the following joke: ``A topologist is a mathematician who cannot tell the difference between his doughnut and coffee cup.'' The joke is a lot funnier\footnote{Eh, maybe.\label{fn-2}} once you've seen \href{https://en.wikipedia.org/wiki/File:Mug_and_Torus_morph.gif}{this animated GIF on Wikipedia}.%
\par
\hypertarget{p-16}{}%
The ``doughnut'''s surface is known mathematically as a ``torus'', shown in \hyperref[figure-torus]{Figure~\ref{figure-torus}}. A sphere is shown in \hyperref[figure-sphere]{Figure~\ref{figure-sphere}}, and a surface that cannot be cannot be embedded in \(\mathbb R^3\), the Klein bottle, is shown in \hyperref[figure-klein-bottle]{Figure~\ref{figure-klein-bottle}}.%
\begin{figure}
\centering
\includegraphics[width=1\linewidth]{images/math.jpg}
\caption{A torus.\label{figure-torus}}
\end{figure}
\begin{figure}
\centering
\includegraphics[width=1\linewidth]{images/math.jpg}
\caption{A sphere.\label{figure-sphere}}
\end{figure}
\begin{figure}
\centering
\includegraphics[width=1\linewidth]{images/math.jpg}
\caption{A Klein bottle.\label{figure-klein-bottle}}
\end{figure}
\hypertarget{p-17}{}%
While these shapes appear very different, they can all be defined as a ``quotient space'' (\hyperref[section-quotient-spaces]{Section~\ref{section-quotient-spaces}}) of the unit square in \(\mathbb R^2\).%
\par
\hypertarget{p-18}{}%
In order to study so-called ``topological spaces'' such as these, we will begin by distilling down the notion of a ``neighborhood'' for an arbitrary set.%
\typeout{************************************************}
\typeout{Section 2 Topological Spaces}
\typeout{************************************************}
\section[{Topological Spaces}]{Topological Spaces}\label{section-topological-spaces}
\typeout{************************************************}
\typeout{Section 3 Continuity \& Homeomorphisms}
\typeout{************************************************}
\section[{Continuity \& Homeomorphisms}]{Continuity \& Homeomorphisms}\label{section-continuity}
\typeout{************************************************}
\typeout{Section 4 Compactness}
\typeout{************************************************}
\section[{Compactness}]{Compactness}\label{section-compactness}
\typeout{************************************************}
\typeout{Section 5 Connectedness}
\typeout{************************************************}
\section[{Connectedness}]{Connectedness}\label{section-connectedness}
\typeout{************************************************}
\typeout{Section 6 Metric Spaces}
\typeout{************************************************}
\section[{Metric Spaces}]{Metric Spaces}\label{section-metric-spaces}
\typeout{************************************************}
\typeout{Section 7 Product Spaces}
\typeout{************************************************}
\section[{Product Spaces}]{Product Spaces}\label{section-product-spaces}
\typeout{************************************************}
\typeout{Section 8 Quotient Spaces}
\typeout{************************************************}
\section[{Quotient Spaces}]{Quotient Spaces}\label{section-quotient-spaces}
\typeout{************************************************}
\typeout{Appendix A Naive Set Theory}
\typeout{************************************************}
%
\appendix
%
\section[{Naive Set Theory}]{Naive Set Theory}\label{set-theory}
\hypertarget{p-19}{}%
A review of basic results concerning sets.%
\end{document}